\documentclass{article}
\usepackage{amsmath,amssymb,amsthm}
\usepackage{algorithm,algorithmic}
\usepackage{bm}
\usepackage{enumitem}
\usepackage{hyperref}
\usepackage{geometry}
\geometry{margin=1in}
\newtheorem{theorem}{Theorem}
\newtheorem{lemma}{Lemma}
\newcommand{\sign}{\operatorname{sign}}
\newcommand{\E}{\mathbb{E}}
\newcommand{\R}{\mathbb{R}}
\newcommand{\norm}[1]{\left\lVert#1\right\rVert}
\newcommand{\inner}[2]{\left\langle #1,#2\right\rangle}
\begin{document}

\section*{Algorithm Name}
\textbf{Momentum Sign Stochastic Gradient Descent (M‑SignSGD)}  

The method updates the parameters using only the sign of a momentum‑averaged stochastic gradient.

\section*{Mathematical Formulation}
Let $f:\R^{d}\rightarrow\R$ be the (possibly non‑convex) objective.  
At iteration $k\ge 0$:

\begin{align}
    g_k &:= \nabla f(x_k;\,\xi_k) \qquad\text{(stochastic gradient, $\xi_k$ random seed)}\label{eq:grad}\\[4pt]
    m_k &:= \beta\, m_{k-1} + (1-\beta)\, g_k ,\qquad m_{-1}=0 \label{eq:momentum}\\[4pt]
    s_k &:= \sign(m_k)\in\{-1,+1\}^{d} \label{eq:sign}\\[4pt]
    x_{k+1} &:= x_k - \alpha_k\, s_k .\label{eq:update}
\end{align}
Here $\beta\in[0,1)$ is the momentum coefficient and $\{\alpha_k\}_{k\ge0}$ is a (possibly decreasing) stepsize sequence.

\section*{Pseudocode}
\begin{algorithm}[H]
\caption{M‑SignSGD}
\begin{algorithmic}[1]
\REQUIRE Initial point $x_0\in\R^{d}$, stepsizes $\{\alpha_k\}$, momentum $\beta\in[0,1)$
\STATE $m_{-1}\gets 0$
\FOR{$k=0,1,\dots,T-1$}
    \STATE Sample a minibatch (or data point) $\xi_k$ and compute stochastic gradient $g_k=\nabla f(x_k;\xi_k)$
    \STATE $m_k \gets \beta\, m_{k-1} + (1-\beta)\, g_k$
    \STATE $s_k \gets \sign(m_k)$ \COMMENT{Componentwise sign}
    \STATE $x_{k+1} \gets x_k - \alpha_k\, s_k$
\ENDFOR
\RETURN $x_T$
\end{algorithmic}
\end{algorithm}

\section*{Dry Run Example}
Consider the one‑dimensional quadratic $f(w)=(w-3)^2$.  
Exact gradient: $\nabla f(w)=2(w-3)$.  
We take deterministic “stochastic’’ gradients (no noise) for illustration.

\begin{center}
\begin{tabular}{c|c|c|c|c}
$k$ & $w_k$ & $g_k=2(w_k-3)$ & $m_k$ & $w_{k+1}=w_k-\alpha\,\sign(m_k)$\\ \hline
0 & $0$ & $-6$ & $m_0=(1-\beta)(-6)$ & $w_1=0-\alpha\cdot(-1)=\alpha$\\[2pt]
1 & $\alpha$ & $2(\alpha-3)$ & $m_1=\beta m_0+(1-\beta)g_1$ & $w_2=w_1-\alpha\,\sign(m_1)$\\[2pt]
2 & $w_2$ & $2(w_2-3)$ & $m_2=\beta m_1+(1-\beta)g_2$ & $w_3=w_2-\alpha\,\sign(m_2)$
\end{tabular}
\end{center}
Take $\alpha=0.1$, $\beta=0.9$:

\begin{itemize}
\item $k=0$: $g_0=-6$, $m_0=-0.6$, $\sign(m_0)=-1$, $w_1=0+0.1=0.1$, $f(w_1)= (0.1-3)^2=8.41$.
\item $k=1$: $g_1=2(0.1-3)=-5.8$, $m_1=0.9(-0.6)+(0.1)(-5.8)=-0.54-0.58=-1.12$, $\sign(m_1)=-1$, $w_2=0.1+0.1=0.2$, $f(w_2)=7.84$.
\item $k=2$: $g_2=2(0.2-3)=-5.6$, $m_2=0.9(-1.12)+(0.1)(-5.6)=-1.008-0.56=-1.568$, $\sign(m_2)=-1$, $w_3=0.2+0.1=0.3$, $f(w_3)=7.29$.
\end{itemize}
The iterates move toward the minimiser $w^\star=3$, decreasing the objective at each step.

\newpage
\section*{Assumptions}
\begin{enumerate}[label={A\arabic*.}, leftmargin=*, itemsep=2pt]
\item \textbf{Smoothness.} $f$ is $L$‑smooth:
\[
\forall x,y\in\R^{d},\qquad 
f(y)\le f(x)+\inner{\nabla f(x)}{y-x}+\frac{L}{2}\norm{y-x}^{2}. \tag{A1}
\]

\item \textbf{Lower Boundedness.} $f^\star:=\inf_{x}f(x)>-\infty$. \tag{A2}

\item \textbf{Unbiased Stochastic Gradient.}
\[
\E[g_k\mid x_k]=\nabla f(x_k),\qquad 
\E\big[\|g_k-\nabla f(x_k)\|^{2}\mid x_k\big]\le\sigma^{2}. \tag{A3}
\]

\item \textbf{Bounded Second Moment.}
\[
\E\big[\|g_k\|^{2}\mid x_k\big]\le G^{2}. \tag{A4}
\]

\item \textbf{Sign Alignment.} There exists a constant $c\in(0,1]$ such that for every $k$
\[
\E\big[\inner{\nabla f(x_k)}{\sign(g_k)}\mid x_k\big]\ge c\,\|\nabla f(x_k)\|_{2}. \tag{A5}
\]
Condition (A5) holds, for example, when each component $g_{k,i}$ is a sub‑Gaussian
random variable with mean $\nabla_i f(x_k)$ and variance bounded by $\sigma^{2}$.
\end{enumerate}

\section*{Main Convergence Theorem}
\begin{theorem}[Convergence of M‑SignSGD]\label{thm:main}
Assume \textbf{A1}–\textbf{A5} and choose a constant stepsize $\alpha_k\equiv\alpha>0$ satisfying
\[
0<\alpha\le \frac{c}{L\sqrt{d}} .
\]
Then after $T$ iterations,
\[
\frac{1}{T}\sum_{k=0}^{T-1}\E\big[\|\nabla f(x_k)\|_{2}\big]
\;\le\;
\frac{f(x_0)-f^\star}{c\alpha T}
\;+\;
\frac{L\alpha d}{2c}.
\tag{1}
\]
In particular, setting $\alpha=\sqrt{\frac{2\,(f(x_0)-f^\star)}{L d\,T}}$ yields
\[
\frac{1}{T}\sum_{k=0}^{T-1}\E\big[\|\nabla f(x_k)\|_{2}\big]
\;=\;O\!\left(\frac{1}{\sqrt{T}}\right).
\]
\end{theorem}

\section*{Theoretical Properties}
\begin{enumerate}[label={P\arabic*.}, leftmargin=*, itemsep=2pt]
\item \textbf{Stability.} Because $\|s_k\|_{2}=\sqrt{d}$ for all $k$, the update magnitude is exactly $\alpha\sqrt{d}$. The stepsize condition $\alpha\le c/(L\sqrt{d})$ guarantees that the quadratic term in the smoothness inequality never outweighs the descent term.

\item \textbf{Hyper‑parameter Sensitivity.}
    \begin{itemize}
        \item Momentum $\beta$ only affects the variance of the sign estimator; the convergence bound does not depend explicitly on $\beta$, but larger $\beta$ typically reduces the noise in $s_k$.
        \item The constant $c$ in (A5) captures how well the sign aligns with the true gradient. Better alignment (larger $c$) improves the rate constant.
    \end{itemize}
\item \textbf{Comparison to Baselines.}
    \begin{itemize}
        \item Standard SGD under the same assumptions yields $\frac{1}{T}\sum\E\|\nabla f(x_k)\|^{2}=O(1/\sqrt{T})$.
        \item M‑SignSGD attains the same $O(1/\sqrt{T})$ rate while communicating only one bit per coordinate per iteration, at the price of a dependence on $\sqrt{d}$ in the stepsize.
    \end{itemize}
\end{enumerate}

\section*{Discussion}
The theorem shows that even though only the sign of a momentum‑averaged stochastic gradient is used, the algorithm enjoys the same sub‑linear convergence guarantee as classical SGD for non‑convex smooth objectives. The proof hinges on the smoothness inequality, boundedness of the sign vector, and the alignment condition (A5) which is satisfied for a broad class of sub‑Gaussian noise models.

\newpage
\section*{Preliminaries (Definitions and Lemmas)}
\begin{lemma}[Smoothness Descent Lemma]\label{lem:smooth}
If $f$ is $L$‑smooth then for any $x$ and any vector $u$,
\[
f(x-u)\le f(x)-\inner{\nabla f(x)}{u}+\frac{L}{2}\|u\|^{2}. 
\tag{L1}
\]
\end{lemma}
\begin{proof}
Apply (A1) with $y=x-u$. \hfill $\blacksquare$
\end{proof}

\begin{lemma}[Sign Alignment Expectation]\label{lem:signalign}
Under (A5) and the momentum update \eqref{eq:momentum},
\[
\E\big[\inner{\nabla f(x_k)}{\sign(m_k)}\mid x_k\big]\ge c\,\|\nabla f(x_k)\|_{2}.
\tag{L2}
\]
\end{lemma}
\begin{proof}
Conditioning on $x_k$, $m_k$ is an affine combination of $g_k$ and $m_{k-1}$. Since $m_{k-1}$ is measurable w.r.t. $\mathcal{F}_{k-1}$, the conditional distribution of $m_k$ given $x_k$ is a shifted version of $g_k$. The sign function is odd, thus
\[
\E[\sign(m_k)\mid x_k]=\E[\sign(g_k)\mid x_k].
\]
Applying (A5) yields (L2). \hfill $\blacksquare$
\end{proof}

\section*{Main Proof (Step‑by‑Step Derivation)}
We prove Theorem~\ref{thm:main}. Let $\mathcal{F}_k$ denote the $\sigma$‑algebra generated by $\{x_0,\dots,x_k\}$.

\begin{enumerate}[label={[\textbf{Step \arabic*}]}, leftmargin=*, itemsep=4pt]
\item \textbf{Apply Smoothness Lemma.}  
Using Lemma~\ref{lem:smooth} with $u=\alpha s_k$,
\[
f(x_{k+1})\le f(x_k)-\alpha\inner{\nabla f(x_k)}{s_k}+\frac{L\alpha^{2}}{2}\|s_k\|^{2}. \tag{2}
\]

\item \textbf{Take Conditional Expectation.}  
Condition on $\mathcal{F}_k$ and use that $s_k=\sign(m_k)$ is $\mathcal{F}_k$‑measurable,
\[
\E\!\left[f(x_{k+1})\mid\mathcal{F}_k\right]
\le f(x_k)-\alpha\,\E\!\left[\inner{\nabla f(x_k)}{s_k}\mid\mathcal{F}_k\right]
     +\frac{L\alpha^{2}}{2}\E\!\left[\|s_k\|^{2}\mid\mathcal{F}_k\right]. \tag{3}
\]

\item \textbf{Bound the Norm of the Sign Vector.}  
Since each component of $s_k$ equals $\pm1$,
\[
\|s_k\|^{2}=d \quad\text{deterministically}. \tag{4}
\]

\item \textbf{Use the Alignment Lemma.}  
From Lemma~\ref{lem:signalign},
\[
\E\!\left[\inner{\nabla f(x_k)}{s_k}\mid\mathcal{F}_k\right]
= \E\!\left[\inner{\nabla f(x_k)}{\sign(m_k)}\mid\mathcal{F}_k\right]
\ge c\,\|\nabla f(x_k)\|_{2}. \tag{5}
\]

\item \textbf{Combine (3)–(5).}  
Plugging (4) and (5) into (3) gives
\[
\E\!\left[f(x_{k+1})\mid\mathcal{F}_k\right]
\le f(x_k)-\alpha c\,\|\nabla f(x_k)\|_{2}
     +\frac{L\alpha^{2}d}{2}. \tag{6}
\]

\item \textbf{Take Full Expectation.}  
Denote $\Phi_k:=\E[f(x_k)]$. Taking expectation of (6) yields
\[
\Phi_{k+1}\le \Phi_k-\alpha c\,\E\!\left[\|\nabla f(x_k)\|_{2}\right]
               +\frac{L\alpha^{2}d}{2}. \tag{7}
\]

\item \textbf{Sum Over Iterations.}  
Summing (7) for $k=0,\dots,T-1$ telescopes the $\Phi$ terms:
\[
\Phi_T\le \Phi_0-\alpha c\sum_{k=0}^{T-1}\E\!\left[\|\nabla f(x_k)\|_{2}\right]
           +\frac{L\alpha^{2}d}{2}T. \tag{8}
\]

\item \textbf{Use Lower Boundedness.}  
Since $\Phi_T\ge f^\star$ by (A2), rearrange (8):
\[
\alpha c\sum_{k=0}^{T-1}\E\!\left[\|\nabla f(x_k)\|_{2}\right]
\le \Phi_0-f^\star+\frac{L\alpha^{2}d}{2}T. \tag{9}
\]

\item \textbf{Divide by $\alpha c T$.}  
\[
\frac{1}{T}\sum_{k=0}^{T-1}\E\!\left[\|\nabla f(x_k)\|_{2}\right]
\le \frac{\Phi_0-f^\star}{c\alpha T}
   +\frac{L\alpha d}{2c}. \tag{10}
\]

\item \textbf{Choose Optimal Constant Stepsize.}  
The right‑hand side of (10) is minimized for
\[
\alpha^{\star}= \sqrt{\frac{2(\Phi_0-f^\star)}{L d\,T}}. \tag{11}
\]
Because of the stepsize restriction $\alpha\le c/(L\sqrt{d})$, the choice (11) satisfies it for all $T\ge1$ (the square‑root term is at most $c/(L\sqrt{d})$ when $T\ge 2(\Phi_0-f^\star)/c^{2}$; otherwise we use $\alpha=c/(L\sqrt{d})$, which yields the same $O(1/\sqrt{T})$ bound up to a constant).

\item \textbf{Insert $\alpha^{\star}$ into (10).}  
Using (11) we obtain
\[
\frac{1}{T}\sum_{k=0}^{T-1}\E\!\left[\|\nabla f(x_k)\|_{2}\right]
\le \sqrt{\frac{L d(\Phi_0-f^\star)}{2}}\,\frac{1}{\sqrt{T}}. \tag{12}
\]
Thus the average expected gradient norm decays as $O(1/\sqrt{T})$, completing the proof. \hfill $\blacksquare$
\end{enumerate}

\section*{Rate Derivation (With Constants)}
From (12) we can read the explicit constant:
\[
\boxed{
\frac{1}{T}\sum_{k=0}^{T-1}\E\!\left[\|\nabla f(x_k)\|_{2}\right]
\le 
\underbrace{\sqrt{\frac{L d(\Phi_0-f^\star)}{2}}}_{\displaystyle C}
\;T^{-1/2}}.
\]
If one prefers a bound on the minimum gradient norm rather than the average, note that
\[
\min_{0\le k<T}\E\!\left[\|\nabla f(x_k)\|_{2}\right]
\le \frac{1}{T}\sum_{k=0}^{T-1}\E\!\left[\|\nabla f(x_k)\|_{2}\right]
\le C\,T^{-1/2}.
\]

\section*{Special Cases}
\begin{enumerate}[label={S\arabic*.}, leftmargin=*, itemsep=2pt]
\item \textbf{Convex $f$.} When $f$ is convex, the same proof yields a bound on the suboptimality gap:
\[
\frac{1}{T}\sum_{k=0}^{T-1}\E\!\big[f(x_k)-f^\star\big]
\le \frac{L d}{2c}\,\alpha +\frac{f(x_0)-f^\star}{c\alpha T},
\]
which after the optimal $\alpha$ choice recovers the $O(1/\sqrt{T})$ rate for convex problems.

\item \textbf{Deterministic Gradients.} If $g_k=\nabla f(x_k)$ (no stochasticity), then $c=1$ and $\sigma=0$. The bound simplifies to
\[
\frac{1}{T}\sum_{k=0}^{T-1}\|\nabla f(x_k)\|_{2}
\le \frac{f(x_0)-f^\star}{\alpha T}+\frac{L\alpha d}{2}.
\]
Choosing $\alpha=\sqrt{2(f(x_0)-f^\star)/(L d T)}$ gives the same $O(1/\sqrt{T})$ rate without any variance term.

\item \textbf{High‑Dimensional Regime.} The dependence on $\sqrt{d}$ appears only in the stepsize constraint and the constant $C$. If communication cost dominates, the sign‑based update remains advantageous because each iteration transmits only $d$ bits irrespective of $L$ or $G$.

\end{enumerate}

\end{document}